\documentclass[11pt]{article}
\usepackage[T1]{fontenc}
\usepackage[utf8]{inputenc}
\usepackage{lmodern,amsmath,amssymb,amsthm,mathtools,microtype}
\usepackage[margin=1in]{geometry}
\usepackage{booktabs,tabularx,xurl}
\usepackage[hidelinks]{hyperref}
\hypersetup{pdftitle={Recognizing Tutte's fourth-kind configuration},pdfauthor={AI-assisted explanatory supplement},pdfsubject={Selected lattice, induced cut, and path recognition}}
\newtheorem{theorem}{Theorem}
\newtheorem{lemma}{Lemma}
\DeclareMathOperator{\rk}{rk}
\DeclareMathOperator{\cl}{cl}
\DeclareMathOperator{\cork}{cork}
\newcommand{\HH}{\mathcal H}
\newcommand{\JJ}{\mathcal J}
\newcommand{\GG}{\Gamma}
\newcommand{\code}[1]{\nolinkurl{#1}}
\setlength{\emergencystretch}{2em}
\title{Recognizing Tutte's fourth-kind configuration}
\author{AI-assisted explanatory supplement to Baker--Jin--Lorscheid}
\date{17 September 2026}
\begin{document}
\maketitle
\begin{abstract}
We give the finite-lattice, modular-cut, and path identification needed to pass from Tutte's six-flat description to the labelled $M(K_{2,3})$ configuration used in [BJL, Remark B.6 and Lemma B.28]. The conclusion concerns the selected join-generated upper sublattice, not necessarily the whole ambient interval. The recognition argument is given in full; a short appendix proves the structural spanning fact it uses from elementary finite-matroid properties. A final section identifies the corresponding existing Lean declarations.
\end{abstract}

\paragraph{Conventions.}
All matroids are finite. The rank and closure in a matroid $N$ are $r_N$ and $\cl_N$. A flat $P$ is \emph{indecomposable} if $N/P$ is connected. We use the rank convention that $N$ is disconnected exactly when
\[
 E(N)=A\sqcup B,\qquad A,B\ne\varnothing,\qquad
 r_N(A)+r_N(B)=r(N).
\]
In particular, empty and one-element matroids are connected. A point is a rank-one flat in a loopless matroid. Joins are ambient closures of unions, and corank means $r(N)-r_N(F)$. A pair of flats is modular when the sum of their ranks equals the sum of the ranks of their meet and join. A modular cut is upward closed among flats and closed under intersections of modular pairs. An upper sublattice is join-closed and geometric in its induced order, with rank equal to the ambient corank of its bottom. Its meets are intrinsic and need not be ambient intersections. A Tutte path is a nonempty finite sequence of hyperplanes with distinct consecutive terms intersecting in indecomposable corank-two flats. It is off a cut if none of its terms is in the cut.

\begin{lemma}\label{lem:two}
A corank-two flat $L$ is decomposable if and only if it is contained in exactly two hyperplanes. Their union is then the whole ground set. If two distinct hyperplanes over $L$ have union equal to the ground set, $L$ is decomposable.
\end{lemma}
\begin{proof}
The contraction $N/L$ is loopless of rank two. If it is disconnected, a nontrivial direct-sum decomposition has two nonempty rank-one summands. Its only hyperplanes are their two ground sets. Lifting gives the assertion. Conversely, a loopless rank-two matroid with precisely two points is their direct sum. If two hyperplanes over $L$ cover the ground set, their images cover $N/L$ by two points; distinct points partition the ground set of a loopless rank-two matroid, so there can be no third one.
\end{proof}

\begin{lemma}\label{lem:span}
If $D\subseteq S$ are indecomposable flats, then $S$ is the join of $D$ and the indecomposable flats $P$ with
\[
 D\subsetneq P\subseteq S,\qquad r(P)=r(D)+1.
\]
Consequently, in a connected loopless matroid every hyperplane is generated by indecomposable points.
\end{lemma}
\noindent A proof is supplied in Appendix~\ref{app:span}. This is the consequence of the Indecomposable Chain and Diamond Properties used in the formalization.

\section{The recognition statement}
\begin{theorem}\label{thm:recognition}
Let $M$ be a finite matroid of rank $r$, let $D$ be an indecomposable flat of corank four, and let $W,Y,T$ be distinct hyperplanes containing $D$. Suppose that
\[
 L_0=W\cap Y,\qquad L_1=W\cap T,\qquad L_2=Y\cap T
\]
are decomposable corank-two flats. Suppose that there are exactly six indecomposable corank-three flats above $D$, and that each $L_i$ contains exactly two of those six flats.

Then the family $\JJ$ consisting of $D$ and all nonempty joins of the six flats is an upper sublattice isomorphic to the flat lattice of $M(K_{2,3})$.

Suppose in addition that $\Gamma$ is a nonempty modular cut, that $W,Y,T\notin\Gamma$, and that each of the six flats is contained in exactly two \emph{hyperplanes} in $\Gamma$. Let
\[
 \delta=(W,X,Y,Z,W)
\]
be a Tutte path off $\Gamma$, such that $W\cap X\cap Y$ and $W\cap Z\cap Y$ are distinct members of the six flats contained in $L_0$. After labelling the pairs by $(1,4),(2,5),(3,6)$, and allowing cyclic rotation of $\delta$, the entire induced cut and the path are
\[
 \Gamma\cap\JJ=
 \{H_{123},H_{156},H_{246},H_{345},E\},
 \qquad
 \delta=(H_{1245},H_{126},H_{1346},H_{456},H_{1245}).
\]
Here $H_S$ means the ambient join of the flats with labels in $S$. The flats with labels $14,25,36$ remain decomposable in $M$. Thus $\delta$ is elementary of the fourth kind in the local [BJL] definition.
\end{theorem}

\paragraph{Relation to the application.}
For Lemma B.28, $D$ is the carrier of the special path, and the preceding counting lemmas supply the six flats and the two-hyperplane incidence counts. The carrier is indecomposable: adjacent hyperplanes have union different from $E$ by Lemma~\ref{lem:two}, and the intersection property proved in Appendix~\ref{app:span} applies successively along the path. The theorem proves the recognition implication used there; it does not assert that an arbitrary local [BJL] configuration forces the same exhaustive ambient hypotheses.

\begin{proof}
Contract $D$ and put $N=M/D$. The contraction is connected and loopless of rank four. Flats above $D$, their joins, hyperplanes, and the cut transport to $N$; indecomposability is preserved by iterated contraction. We use the same letters for the transported objects. The six distinguished flats are now points, and the $L_i$ are lines.

\paragraph{1. The six points split into three disjoint pairs.}
By Lemma~\ref{lem:two},
\[
 W\cup Y=W\cup T=Y\cup T=E(N).
\]
Every element therefore belongs to at least two of $W,Y,T$. If $P$ is a point, choose $e\in P$; since $P=\cl_N(e)$, it lies in one of $L_0,L_1,L_2$. The lines are distinct: an equality between two would put a decomposable corank-two flat under three distinct hyperplanes.

Every distinguished point belongs to at least one of these lines. There are six such points, and the sum of the three specified counts is $2+2+2=6$. Thus each belongs to exactly one line. Write the pairs on $L_0,L_1,L_2$ as $P,Q,R$. The two distinct points in each pair generate their line. Their incidences with the hyperplanes are precisely
\[
 W:\ P\cup Q,\qquad Y:\ P\cup R,\qquad T:\ Q\cup R.
\]
In particular, a point in one pair is outside the hyperplane containing the other two pairs.

\paragraph{2. The selected restriction is $M(K_{2,3})$.}
Choose a representative element of each distinguished point. Every triple of representatives is independent. A triple containing a complete pair has rank three because its third point is outside that pair's line. A triple taking one point from each pair has rank three because the first two lie in their pair-hyperplane and the third is outside it.

A four-set consisting of two complete pairs has rank three. Every other four-set has distribution $(2,1,1)$ among the pairs. It spans $N$: a hypothetical containing hyperplane would contain the line generated by the complete pair, so Lemma~\ref{lem:two} would make it one of that line's two hyperplanes. Neither contains points from both remaining pairs. This is impossible.

Label the pairs $P=(1,4)$, $Q=(2,5)$, $R=(3,6)$. The selected restriction has rank four, all triples are independent, and its only dependent four-sets are
\[
 1245,\qquad1346,\qquad2356.
\]
This is exactly the cycle matroid of $K_{2,3}$ when its three two-edge paths have those pairs: its four-edge cycles are the unions of two such paths, and all other four-edge subsets are trees. The larger subsets are determined by their four-subsets.

For completeness, the selected flat lattice has bottom, six singleton flats, fifteen two-element flats, eight transversal triples (one label from each pair), the three displayed four-element hyperplanes, and top: $34$ flats in total. The closure map from this restriction to $N$ is injective on flats, because intersecting an ambient closure with the selected representatives recovers the restriction flat. It preserves order, joins, and ranks. The representatives span $N$, so its image has top $E(N)$ and rank four. This image is the required upper sublattice. Lifting back to $M$ gives $\JJ$ with bottom $D$.

\paragraph{3. The hyperplanes and the induced cut.}
Every ambient hyperplane of $N$ is generated by indecomposable points by Lemma~\ref{lem:span}. By the exhaustion assumption, these are among the selected six. Thus every ambient hyperplane is represented by the selected flat lattice: it is $W,Y,T$ or one of the eight transversal-triple hyperplanes. This conclusion about hyperplanes does not exclude additional \emph{decomposable points} or other extra flats in the ambient lattice.

Encode a transversal triple by $abc\in\{0,1\}^3$, choosing one of the two points in each pair. Only these triple-hyperplanes can be in $\Gamma$. Each contains exactly three distinguished points, and each point lies in exactly two cut hyperplanes. Counting incidences gives
\[
 3\,\#\{\hbox{cut hyperplanes}\}=6\cdot2,
\]
so exactly four of the eight triple-hyperplanes belong to the cut.

Two cut triples cannot agree in two coordinates. Otherwise their distinct hyperplanes intersect in the line generated by the two common points: the intersection contains that rank-two flat, and, being strictly below a hyperplane, has rank at most two. The hyperplanes therefore form a modular pair. Their line lies in one of $W,Y,T$, so modular-cut closure and upward closure would put that off-cut hyperplane in $\Gamma$.

The four selected binary triples consequently form one parity class. To see this without a case table, relabel within pairs so that $000$ is selected. Its three neighbors are excluded. If $111$ were selected, its three neighbors would also be excluded, leaving only two possible selected vertices. Hence the remaining selected triples are $011,101,110$. These are exactly
\[
 123,\qquad156,\qquad246,\qquad345.
\]
Every selected flat of rank at most two is contained in one of $W,Y,T$, so none is in the cut. The top is in the cut by nonemptiness. This proves equality of the \emph{entire} induced cut, not just agreement on its hyperplanes.

\paragraph{4. The original path has the prescribed word.}
The middle hyperplanes $X,Z$ are not $W,Y$, since $\delta$ is a Tutte path. Neither is $T$: the intersection $W\cap Y\cap T$ contains no distinguished point. Thus they are transversal-triple hyperplanes. Their triples choose different points of $P$, by the hypothesis about $W\cap X\cap Y$ and $W\cap Z\cap Y$.

Both triples have the parity opposite to the cut. Their Hamming distance is therefore two, not one or three. Since they differ in the $P$ coordinate, they share exactly one of the $Q,R$ coordinates. If necessary, interchange $W,Y$ and cyclically rotate the path by two vertices. This interchanges $Q,R$ and allows the shared coordinate to be $R$.

Label the points of $P,Q$ chosen by $X$ as $1,2$, the other points as $4,5$, and their shared $R$ point as $6$ (the other is $3$). Then $X=H_{126}$ and $Z=H_{456}$. Both have odd parity in these labels, so the cut is the even class displayed above. Moreover $W=H_{1245}$ and $Y=H_{1346}$. This gives exactly the required word and cut simultaneously. The three pair-flats are the original $L_i$, so their ambient decomposability is unchanged.
\end{proof}

\appendix
\section{Proof of the structural spanning lemma}\label{app:span}
We give the elementary details to avoid importing the homotopy theorem, the Special Lemma, or the desired recognition as background. We use only finite-matroid rank submodularity, basis extension, the contraction-rank formula, and the fact that a flat avoiding an element extends to a hyperplane avoiding that element.

\paragraph{A rank-partition observation.}
If $E(N)=A\sqcup B$ and $r_N(A)+r_N(B)=r(N)$, then for every $S\subseteq E(N)$,
\[
 r_N(S)=r_N(S\cap A)+r_N(S\cap B).
\]
Indeed, add the submodular inequalities for $(S,A)$ and $(S\cup A,B)$; subadditivity gives the reverse inequality. The contraction-rank formula then gives, for every $G\subseteq E(N)$,
\[
 r_{N/G}(A-G)+r_{N/G}(B-G)=r(N)-r_N(G)=r(N/G).
\]
If $N/G$ is connected, one of $A,B$ must therefore be contained in $G$.

Conversely, if $A,B$ are a nonempty partition and every hyperplane contains $A$ or $B$, choose a basis $I$ of $A$, extend it to a basis $J$ of $N$, and put $K=J\cap B$. Then $J\cap A=I$. If $K$ did not span $B$, a hyperplane containing $\cl_N(K)$ and avoiding some $b\in B$ would have to contain $A$, hence the basis $J$, a contradiction. Thus $r_N(A)+r_N(B)=|J|=r(N)$.

\paragraph{The intersection property.}
If $F_1,F_2$ are indecomposable and $F_1\cup F_2\ne E$, then $F=F_1\cap F_2$ is indecomposable. For otherwise a nontrivial rank-additive partition $A\sqcup B$ of $M/F$ persists after contracting either $G_i=F_i-F$. As $M/F_i$ is connected, each $G_i$ contains an entire side. They cannot contain the same nonempty side because $G_1\cap G_2=\varnothing$; containing opposite sides forces $F_1\cup F_2=E$. This is a contradiction.

\paragraph{The step property.}
For indecomposable $D\subsetneq S$, there is an indecomposable $U$ between them of rank $r(S)-1$. If $S=E$, choose any hyperplane above $D$. Otherwise the nonempty sets $S-D$ and $E-S$ partition the ground of the connected matroid $M/D$. By the preceding observation there is a hyperplane $H\supseteq D$ with $S\not\subseteq H$ and $S\cup H\ne E$. Choose one maximizing $|S\cap H|$ and put $U=S\cap H$. The intersection property makes $U$ indecomposable.

If $r(U)\leq r(S)-2$, choose $a\in S-H$ and put $U'=\cl(U\cup\{a\})\subsetneq S$. Choose a hyperplane $H'\supseteq U'$ avoiding an element of $S-U'$. If $S\cup H'=E$, then $H'\supseteq E-S$ and $H'\supseteq H\cap S$, hence $H\subseteq H'$, forcing equality; but $a\in H'-H$. Thus $H'$ is another admissible hyperplane and $|H'\cap S|>|H\cap S|$, a contradiction. This proves the step property and, by repetition, the chain property.

\paragraph{The diamond property.}
Suppose $D\subseteq S$ are indecomposable and $r(S)=r(D)+2$. Choose an indecomposable $U$ of intermediate rank and $a\in S-U$; put $V_0=\cl(D\cup\{a\})$. If $V_0$ is indecomposable, we are done. Otherwise $S\ne E$, since $V_0$ would be a hyperplane if $S=E$. Choose a flat $L\supseteq D$ of corank two with $L\cap S=D$ and $L\vee S=E$. Such an $L$ is obtained in $M/D$ by extending a basis of $S-D$ to a basis of the ground set and taking the closure of the added basis elements. Rank submodularity proves the asserted intersection.

Put $H_1=U\vee L$ and $H_2=V_0\vee L$. These are distinct hyperplanes with $S\cap H_1=U$, $S\cap H_2=V_0$, and $H_1\cap H_2=L$. The intersection property gives $S\cup H_2=E$. We have $U\cup H_2\ne E$: otherwise $(U-D)\sqcup(H_2-D)$ would be a nontrivial rank-additive partition of $M/D$. Similarly $S\cup H_1\ne E$, using connectedness of $M/U$ and the partition induced by $S,H_1$.

Choose $b\notin U\cup H_2$ and $c\notin S\cup H_1$, and set $V=\cl(D\cup\{b\})$, $H=V\vee L$. Then $b\in S$, $V\ne U$, $H$ is a hyperplane, and $S\cap H=V$, by the same rank calculations. Since $b\in H-H_2$, we have $H\cap H_2=L$. Also $c\in H_2$; if $c\in H$, then $c\in L\subseteq H_1$, a contradiction. Thus $S\cup H\ne E$, and the intersection property makes $V$ indecomposable. This proves the diamond property.

\paragraph{Generation.}
Induct on $r(S)-r(D)$. The cases zero and one are immediate. Otherwise choose an indecomposable flat $R$ of rank $r(S)-2$ between $D$ and $S$ by the chain property. The diamond property gives distinct indecomposable $U,V$ of rank $r(S)-1$ between $R$ and $S$. They join to $S$, and the inductive assertion for each proves Lemma~\ref{lem:span}. A hyperplane is indecomposable because its contraction is loopless of rank one, so the final assertion of that lemma follows as well.

\clearpage
\section{Provenance and correspondence with Lean}
This is an AI-assisted explanatory note prepared using ChatGPT (GPT-6 Astra Pro). It is separate from the authors' Appendix B. Its inclusion in the repository does not assert independent human review of this prose. Lean checks formal declarations, not a PDF; the correspondence described here is a mathematical explanation of existing code, not a new Lean compilation of this note.

The relevant proof snapshot is public commit
\href{https://github.com/mbaker386/tutte-homotopy-lean/tree/2ff7110d59df7b75e5d00db06dee6ca674e093e6}{\texttt{2ff7110d59df7b75e5d00db06dee6ca674e093e6}}.
All paths below are relative to \code{TutteFormalization/Homotopy/}; declaration names are in \code{TutteFormalization.Homotopy}, with the indicated subnamespaces.

\begin{center}
\renewcommand{\arraystretch}{1.2}
\begin{tabularx}{\textwidth}{@{}p{0.26\textwidth}X@{}}
\toprule
Mathematical step & Formal location and declaration \\\midrule
Spanning lemma & \code{PointGeneration.lean}: \code{indecomposable_subset_of_covers_subset} \\
Selected rank table and embedding & \code{PairPlaneRecognition.lean}: \code{fourthModelOfPairPlanes}; \code{RecognizedSixModel.lean}: \code{SixFlatData.ambientModel} \\
All ambient hyperplanes & \code{HyperplaneRecognition.lean}: \code{fourth_hyperplanes_recognized} \\
Four cut hyperplanes and parity & \code{SixModelCubes.lean}: \code{SixFlatData.cutTriples_card_four}; \code{CutParityRecognition.lean}: \code{actual_cut_cube_parity} \\
Entire cut and original word & \code{FourthCutRecognition.lean}: \code{exactFourthCut_of_hyperplanes}; \code{FourthWordRelabelling.lean}: \code{fourth_word_normalization} \\
Recognition in B.28 & \code{FourthRecognition.lean}: \code{fourth_elementary_of_count}; \code{FourthSpecial.lean}: \code{SixFlatData.original_path_elementary} \\
\bottomrule
\end{tabularx}
\end{center}

The last declaration supplies the recognition needed in the actual Special Lemma proof. Its lower-corank and non-nullness hypotheses are used in deriving the counting data; Theorem~\ref{thm:recognition} starts from those explicit geometric and incidence data. We do not claim a separate single Lean declaration with every hypothesis of this note verbatim. The auxiliary rank-partition, step, and diamond arguments above explain the independently proved infrastructure used by the formal spanning lemma.

\paragraph{Reference.}
[BJL] Matthew Baker, Tong Jin, and Oliver Lorscheid, \emph{A modern perspective on Tutte's homotopy theorem}, with Appendix B by Ju\v s Kocutar, revised manuscript (2026), Section 1.4 and Appendices A--B. This note addresses the comparison in Remark B.6 and its use in Lemma B.28. No conclusion of the homotopy theorem is used to prove the recognition statement.
\end{document}
